\section{Discussion and limitations}
\label{sec:discussion}

The result establishes a nontrivial confined solution with authenticated local
uniqueness and several physical diagnostics.  It does not establish that the
chosen boundary condition is dynamically generated.  The imposed Dirichlet
condition fixes topological charge, represents external support, enters the
radial fluctuation domain, and determines the leading ultraviolet phase and
coefficient.  The support traction and its dynamics are not derived.  The
boundary condition is therefore part of the mathematical model rather than a
numerical convenience.

This distinction matters for smoothing.  A generic finite-stiffness Robin
condition is not automatically a nearby realization in the same exact
topological sector: imposing both the exact boundary value needed for unit
charge and the associated finite-stiffness condition can force vanishing value
and derivative, hence the trivial solution by ODE uniqueness.  A controlled
smooth-confinement limit would require explicit boundary degrees of freedom or
a topology-preserving matter-plus-boundary action.  Standard Fourier endpoint
theory then predicts how the tail changes, but such a family is not constructed
here.

The fixed-wall radial gap is similarly specific.  It excludes an unstable
spherically symmetric profile fluctuation on the declared form domain.  It
does not control a moving membrane, nonspherical grand-spin channels, an anchor
or holding apparatus, or the coupled metric.  No claim of full dynamical
stability follows.

Finally, the metric is prescribed.  We do not solve the Einstein--Skyrme
constraints, impose Israel matching conditions, or prove weak gravitational
backreaction.  The rotor inertia and its finite optical transform make the
profile suitable as an input to those questions, but they do not answer them.
Turning that transform into a bath factor would require a specified current,
interaction, switching protocol, and open-system limit.

The analytic continuation result in \cref{cor:local-parameter-branch} proves
that the solution is not isolated in parameter space, but it supplies no
numerical width.  The natural next mathematical step is an explicit rational
parameter box, with the same certificate architecture and a carefully modeled
boundary.
The natural physical step is a local matter-plus-support action whose static
solution reduces to the validated branch and whose radial, nonspherical, and
gravitational perturbations can be controlled together.  Those extensions
would be new theorems, not consequences silently imported into the present
one.
